% \section{Conclusion}

% We introduced a symbolic dynamic programming framework for algorithms on tree
% decompositions, formalized through Symbolic DP-cores. We
% showed that several natural graph classes admit compact symbolic representations
% theoretically, and that, more generally, symbolic representations can remain
% compact even when treewidth is large, yielding substantial practical gains over
% standard DP. As directions for future work, it would be natural
% to develop SDP-cores for a broader range of graph problems and to identify
% additional graph classes admitting low SDP-width under suitable decompositions.
% Another promising avenue is the optimization of symbolic DP implementations to
% better compete in running time, particularly on dense and highly symmetric graph
% classes where treewidth alone provides limited guidance.

\section{Conclusion}

We presented a symbolic dynamic programming framework for tree decompositions in
which witness sets are represented implicitly using OBDDs. Instantiated for
{\sc Vertex Cover}, the resulting algorithm runs in time
$2^{\BigO(k)}\cdot n^{\BigO(1)}$ on a width-$k$ decomposition, matching the
expected single-exponential dependence on the decomposition width.

Our experiments indicate that, symbolic representations can offer meaningful practical gains over explicit witness tables. On random instances, OBDD-based witness sets remained orders of magnitude
smaller than their explicit counterparts and enabled processing substantially
larger decomposition widths than standard DP. On rook graphs symbolic DP remained tractable well beyond the point where explicit
DP exhausted memory. These results suggest that symbolic representations can
extend the practical reach of decomposition-based dynamic programming on
instances where treewidth alone offers limited guidance.